\thispagestyle{otFront}\phantomsection\pdfbookmark[0]{Avant-propos}{bm-avant-propos}
\otPageTitre{Avant-propos}{Comprendre, chercher et progresser pendant toute la Seconde}
{\large\itshape\color{otNavy}Ce livre accompagne ton année de Seconde générale et technologique, depuis les applications directes jusqu’aux problèmes qui demandent de prendre des initiatives.}\par\bigskip
Les 16 chapitres couvrent les nombres et le calcul, les fonctions, la géométrie, les statistiques et les probabilités. Python, la logique et les automatismes se travaillent tout au long de l’année : les chapitres 15 et 16 et quatre séries de questions flash corrigées servent de ressources à reprendre au moment utile.\par\medskip
Chaque chapitre propose un cours de référence, des méthodes et 15 exercices en trois étapes : \textbf{Découverte}, \textbf{Entraînement}, \textbf{Problèmes}. Les corrigés expliquent pourquoi la méthode convient, déroulent les calculs et interprètent les résultats. Les figures utiles restent avec le texte correspondant.\par\medskip
Un parcours de \textbf{16 défis facultatifs}, avec deux indications progressives pour chacun, permet d’approfondir. Il traite toutes les pistes d’approfondissement citées dans le programme, dont les barycentres, les concours de droites, les identités et les lieux géométriques. Les notions supplémentaires sont définies sur place.\par\medskip
Les démonstrations de référence, deux bilans de fin d’année, un atelier numérique et une matrice de couverture du programme complètent le livre. Les bilans sont des entraînements éditoriaux : ils ne représentent pas une épreuve nationale de Seconde.\par\medskip
Prends le temps de chercher avant de consulter une solution. Après lecture, ferme le corrigé et refais l’exercice : expliquer la démarche avec tes propres mots permet de mesurer ce qui est acquis.\par\bigskip
{\otMedium Bon travail et bonne année de Seconde !}\par\medskip
\includegraphics[width=10mm]{icone-couleur}\clearpage